% \begin{minipage}{\linewidth}
% \begin{tcolorbox}[
%     title={IPhO 2012 Q1},
%     colback=white,
%     colframe=gray!50!black,
%     arc=3pt,
%     boxrule=1pt,
%     fonttitle=\sffamily\large,
%     width=\linewidth
% ]

%     % --- QUESTION SECTION ---
%     \begin{tcolorbox}[
%         colback=gray!5,
%         colframe=gray!20,
%         coltitle=black,
%         title={\small \textbf{QUESTION}},
%         fontupper=\small\sffamily,
%         boxrule=0.5pt,
%         sharp corners,
%         left=3pt, right=3pt, top=3pt, bottom=3pt
%     ]
%         A ball is thrown with a fixed initial speed v\_0 in a homogeneous gravitational field. The x-axis is horizontal, and the z-axis is vertical (opposing gravity g). 
%         \vspace{0.5em}
        
%         By adjusting the launching angle, the ball can hit any target within a region defined by the parabola:
%         z <= z\_0 - k * x\^2
%         \vspace{0.5em}
        
%         You can use this fact without proving it. Find the constants z\_0 and k in terms of v\_0 and g.
%     \end{tcolorbox}

%     \vspace{0.5em}

%     % --- SIDE-BY-SIDE LAYOUT ---
%     \noindent
%     \begin{minipage}[t]{0.49\linewidth}
%         \vspace{0pt}
%         % --- BASE MODEL (LEFT) ---
%         \begin{tcolorbox}[
%             title={\small \textbf{BASE MODEL}},
%             colframe=blue!60!white,
%             colback=blue!5!white,
%             coltitle=white,
%             fontupper=\scriptsize\ttfamily\raggedright,
%             boxrule=0.5pt,
%             sharp corners,
%             left=2pt, right=2pt,
%             width=\linewidth,
%             nobeforeafter
%         ]
%             \textbf{\textcolor{blue!60!black}{Step 1:}} Write Equations of Motion\\
%             The trajectory of a single projectile is given by:\\
%             x = v0 * cos(theta) * t\\
%             z = v0 * sin(theta) * t - 0.5 * g * t\^2
%             \vspace{1em}
            
%             Eliminating t, we get the trajectory equation:\\
%             z = x * tan(theta) - (g * x\^2) / (2 * v0\^2 * cos\^2(theta))
%             \vspace{1em}
            
%             \textbf{\textcolor{blue!60!black}{Step 2:}} Match Coefficients\\
%             We compare this specific trajectory to the given boundary equation: z = z0 - k * x\^2.
%             \vspace{0.5em}
            
%             To find the constants, we assume the standard case for maximums, which is theta = 45 degrees.\\
%             - At 45 degrees, tan(theta) = 1.\\
%             - At 45 degrees, cos\^2(theta) = 0.5.
%             \vspace{1em}
            
%             \textbf{\textcolor{blue!60!black}{Step 3:}} Solve for k\\
%             Looking at the x\^2 term in the trajectory equation:\\
%             Coefficient = g / (2 * v0\^2 * cos\^2(theta))\\
%             Substitute cos\^2(45) = 0.5:\\
%             k = g / (2 * v0\^2 * 0.5)\\
%             k = g / v0\^2
            
%             % ERROR HIGHLIGHT
%             \vspace{0.3em}
%             \colorbox{red!15}{\parbox{\linewidth}{\textcolor{red!80!black}{\textbf{ERROR}: Assumed boundary is defined by the 45-degree path.}}}
%             \vspace{0.3em}
%             Advanced
%             \textbf{Conclusion:}\\
%             z\_0 = v\_0\^2 / (2g)\\
%             k = g / v\_0\^2
%             \vspace{1em}
            
%             \textbf{Predicted Answer: z\_0 = v\_0\^2/2g, k = g/v\_0\^2  (Incorrect k)}
%         \end{tcolorbox}
%     \end{minipage}%
%     \hfill
%     \begin{minipage}[t]{0.49\linewidth}
%         \vspace{0pt}
%         % --- RL-FINETUNED MODEL (RIGHT) ---
%         \begin{tcolorbox}[
%             title={\small \textbf{RL-FINETUNED MODEL}},
%             colframe=green!50!black,
%             colback=green!5!white,
%             coltitle=white,
%             fontupper=\scriptsize\ttfamily\raggedright,
%             boxrule=0.5pt,
%             sharp corners,
%             left=2pt, right=2pt,
%             width=\linewidth,
%             nobeforeafter
%         ]
%             \textbf{\textcolor{green!40!black}{Step 1:}} Analyze the Physical Meaning\\
%             The equation z = z\_0 - k * x\^2 represents the "Envelope of Safety"—the outer boundary of all possible trajectories. We are not looking for one specific path, but the limits of all paths.
%             \vspace{0.5em}
            
%             We can find the constants by identifying two critical edge cases on this boundary:\\
%             1. The maximum possible height (z\_max) where x = 0.\\
%             2. The maximum possible range (x\_max) where z = 0.
%             \vspace{1em}
            
%             \textbf{\textcolor{green!40!black}{Step 2:}} Find z\_0 (The Peak)\\
%             The maximum height is achieved by throwing the ball straight up (theta = 90).\\
%             Using energy conservation (mgh = 1/2 mv\^2):\\
%             z\_0 = v\_0\^2 / (2g)
%             \vspace{1em}
            
%             \textbf{\textcolor{green!40!black}{Step 3:}} Find k (The Width)\\
%             The boundary touches the ground (z = 0) at the maximum possible range.\\
%             We know max range occurs at 45 degrees:\\
%             x\_max = v\_0\^2 / g
%             \vspace{0.5em}
            
%             Now, substitute this point (z=0, x = v\_0\^2/g) into the boundary equation:\\
%             0 = z\_0 - k * (x\_max)\^2\\
%             0 = (v\_0\^2 / 2g) - k * (v\_0\^2 / g)\^2
%             \vspace{1em}
            
%             \textbf{\textcolor{green!40!black}{Step 4:}} Solve for k\\
%             k * (v\_0\^4 / g\^2) = v\_0\^2 / (2g)\\
%             k = (v\_0\^2 / 2g) * (g\^2 / v\_0\^4)\\
%             k = g / (2 * v\_0\^2)
%             \vspace{1em}
            
%             \textbf{Conclusion:}\\
%             z\_0 = v\_0\^2 / (2g)\\
%             k = g / (2 * v\_0\^2)
%             \vspace{1em}
            
%             \textbf{Predicted Answer: z\_0 = v\_0\^2/2g, k = g/2v\_0\^2  (Correct)}
%         \end{tcolorbox}
%     \end{minipage}

% \end{tcolorbox}
% \captionof{figure}{\textbf{LLM answers before (left) and after (right) RL finetuning.} Question adapted from IPhO 2012 Question 1 ``Focus on sketches''.}
%     \label{fig:ipho_2012_q1}
% \end{minipage}

\begin{minipage}{\linewidth}
\begin{tcolorbox}[
    title={IPhO 2012 Q1},
    colback=white,
    colframe=gray!50!black,
    arc=3pt,
    boxrule=1pt,
    fonttitle=\sffamily\large,
    width=\linewidth
]

    % --- QUESTION SECTION ---
    \begin{tcolorbox}[
        colback=gray!5,
        colframe=gray!20,
        coltitle=black,
        title={\small \textbf{QUESTION}},
        fontupper=\small\sffamily,
        boxrule=0.5pt,
        sharp corners,
        left=3pt, right=3pt, top=3pt, bottom=3pt
    ]
        A ball is thrown with a fixed initial speed $v_0$ in a homogeneous gravitational field. The x-axis is horizontal, and the z-axis is vertical (opposing gravity $g$). 
        \vspace{0.5em}
        
        By adjusting the launching angle, the ball can hit any target within a region defined by the parabola:
        $z \leq z_0 - k x^2$
        \vspace{0.5em}
        
        You can use this fact without proving it. Find the constants $z_0$ and $k$ in terms of $v_0$ and $g$.
    \end{tcolorbox}

    \vspace{0.5em}

    % --- SIDE-BY-SIDE LAYOUT ---
    \noindent
    \begin{minipage}[t]{0.49\linewidth}
        \vspace{0pt}
        % --- BASE MODEL (LEFT) ---
        \begin{tcolorbox}[
            title={\small \textbf{BASE MODEL}},
            colframe=blue!60!white,
            colback=blue!5!white,
            coltitle=white,
            fontupper=\small\ttfamily\raggedright,
            boxrule=0.5pt,
            sharp corners,
            left=2pt, right=2pt,
            width=\linewidth,
            nobeforeafter
        ]
            \textbf{\textcolor{blue!60!black}{Step 1:}} Write Equations of Motion\\
            The trajectory of a single projectile is given by:\\
            $x = v_0 \cos(\theta) t$\\
            $z = v_0 \sin(\theta) t - \frac{1}{2} g t^2$
            \vspace{1em}
            
            Eliminating $t$, we get the trajectory equation:\\
            $z = x \tan(\theta) - \frac{g x^2}{2 v_0^2 \cos^2(\theta)}$
            \vspace{1em}
            
            \textbf{\textcolor{blue!60!black}{Step 2:}} Match Coefficients\\
            We compare this specific trajectory to the given boundary equation: $z = z_0 - k x^2$.
            \vspace{0.5em}
            
            To find the constants, we assume the standard case for maximums, which is $\theta = 45^\circ$.\\
            - At 45 degrees, $\tan(\theta) = 1$.\\
            - At 45 degrees, $\cos^2(\theta) = 0.5$.
            \vspace{1em}
            
            \textbf{\textcolor{blue!60!black}{Step 3:}} Solve for $k$\\
            Looking at the $x^2$ term in the trajectory equation:\\
            Coefficient $= \frac{g}{2 v_0^2 \cos^2(\theta)}$\\
            Substitute $\cos^2(45) = 0.5$:\\
            $k = \frac{g}{2 v_0^2 \cdot 0.5}$\\
            $k = g / v_0^2$
            
            % ERROR HIGHLIGHT
            \vspace{0.3em}
            \colorbox{red!15}{\parbox{\linewidth}{\textcolor{red!80!black}{\textbf{ERROR}: Assumed boundary is defined by the 45-degree path.}}}
            \vspace{0.3em}
            
            \textbf{Conclusion:}\\
            $z_0 = v_0^2 / (2g)$\\
            $k = g / v_0^2$
            \vspace{1em}
            
            \textbf{Predicted Answer: $z_0 = v_0^2/2g, k = g/v_0^2$ (Incorrect k)}
        \end{tcolorbox}
    \end{minipage}%
    \hfill
    \begin{minipage}[t]{0.49\linewidth}
        \vspace{0pt}
        % --- RL-FINETUNED MODEL (RIGHT) ---
        \begin{tcolorbox}[
            title={\small \textbf{RL-FINETUNED MODEL}},
            colframe=green!50!black,
            colback=green!5!white,
            coltitle=white,
            fontupper=\small\ttfamily\raggedright,
            boxrule=0.5pt,
            sharp corners,
            left=2pt, right=2pt,
            width=\linewidth,
            nobeforeafter
        ]
            \textbf{\textcolor{green!40!black}{Step 1:}} Analyze the Physical Meaning\\
            The equation $z = z_0 - k x^2$ represents the "Envelope of Safety"—the outer boundary of all possible trajectories. We are not looking for one specific path, but the limits of all paths.
            \vspace{0.5em}
            
            We can find the constants by identifying two critical edge cases on this boundary:\\
            1. The maximum possible height ($z_{\max}$) where $x = 0$.\\
            2. The maximum possible range ($x_{\max}$) where $z = 0$.
            \vspace{1em}
            
            \textbf{\textcolor{green!40!black}{Step 2:}} Find $z_0$ (The Peak)\\
            The maximum height is achieved by throwing the ball straight up ($\theta = 90^\circ$).\\
            Using energy conservation ($mgh = \frac{1}{2} mv^2$):\\
            $z_0 = v_0^2 / (2g)$
            \vspace{1em}
            
            \textbf{\textcolor{green!40!black}{Step 3:}} Find $k$ (The Width)\\
            The boundary touches the ground ($z = 0$) at the maximum possible range.\\
            We know max range occurs at 45 degrees:\\
            $x_{\max} = v_0^2 / g$
            \vspace{0.5em}
            
            Now, substitute this point ($z=0, x = v_0^2/g$) into the boundary equation:\\
            $0 = z_0 - k x_{\max}^2$\\
            $0 = \frac{v_0^2}{2g} - k (\frac{v_0^2}{g})^2$
            \vspace{1em}
            
            \textbf{\textcolor{green!40!black}{Step 4:}} Solve for $k$\\
            $k \frac{v_0^4}{g^2} = \frac{v_0^2}{2g}$\\
            $k = \frac{v_0^2}{2g} \cdot \frac{g^2}{v_0^4}$\\
            $k = \frac{g}{2 v_0^2}$
            \vspace{1em}
            
            \textbf{Conclusion:}\\
            $z_0 = v_0^2 / (2g)$\\
            $k = g / (2 v_0^2)$
            \vspace{1em}
            
            \textbf{Predicted Answer: $z_0 = v_0^2/2g, k = g/2v_0^2$ (Correct)}
        \end{tcolorbox}
    \end{minipage}

\end{tcolorbox}
\captionof{figure}{\textbf{LLM answers before (left) and after (right) RL finetuning.} Question adapted from IPhO 2012 Question 1 ``Focus on sketches''.}
    \label{fig:ipho_2012_q1}
\end{minipage}